\chapter{Occupational Health and Dust Risk Assessment}
\label{chap:Occupational Health and Dust Risk Assessment}

%---------------------------- SECTION ----------------------------
\section{Why this assessment matters in construction}

\paragraph{Occupational disease kills more construction workers in the United Kingdom each year than all acute traumatic injuries combined. The HSE's annual statistics on work-related deaths consistently reveal that the 35 to 60 fatal accidents in construction each year are dwarfed by an estimated 3,500 to 5,000 construction workers who die annually from past occupational exposures --- principally mesothelioma from asbestos, chronic obstructive pulmonary disease (COPD) and silicosis from respirable dust, and lung cancer from silica and diesel engine emissions. These deaths do not appear in the year's incident statistics because the disease that causes them was contracted decades earlier, and the latency between exposure and diagnosis may be 10 to 40 years. They are invisible until they kill, and the invisibility of the mechanism --- a worker coughing more than usual, feeling breathless on the stairs, noticing declining stamina --- means that the hazard receives far less management attention than a fall from height or a struck-by incident, despite its far greater mortality impact.}

\paragraph{Respirable crystalline silica (RCS) is the single most significant occupational health hazard in the construction industry after asbestos. Silica is present in virtually every construction material: concrete, mortar, brick, block, sandstone, granite, and engineered stone all contain crystalline silica at levels that, when disturbed by cutting, grinding, drilling, or breaking, generate respirable dust particles that penetrate to the alveolar region of the lung and cause irreversible fibrotic scarring (silicosis). Silicosis is incurable and progressive; it may also predispose to tuberculosis and lung cancer. The Workplace Exposure Limit (WEL) for RCS under COSHH Regulations 2002 is 0.1 mg/m\textsuperscript{3} as an 8-hour time-weighted average --- a threshold that is routinely exceeded on construction sites where dry cutting, dry grinding, and dry drilling on silica-containing materials are conducted without engineering controls. The HSE's campaign on silica dust (``Dust kills --- control it now'') reflects the evidential basis that uncontrolled silica exposure on construction sites is causing an ongoing epidemic of silicosis and silica-related lung cancer that will not become statistically visible for another 20--30 years.}

\paragraph{Wood dust, including hardwood dust and, to a lesser extent, softwood dust, is a classified carcinogen causing sinonasal cancer and nasal adenocarcinoma. The WEL for hardwood dust is 3 mg/m\textsuperscript{3} and for softwood dust 5 mg/m\textsuperscript{3} as an 8-hour TWA, but recent HSE guidance suggests that these limits require review given emerging evidence of cancer risk at lower exposures. Construction joiners, site carpenters, and wood machining operatives are the principal risk groups. Diesel engine emissions from construction plant operating in confined or poorly ventilated spaces generate an aerosol mixture of particulates and gases, including diesel particulate matter (DPM), which is classified as a Group~1 carcinogen by the International Agency for Research on Cancer (IARC) and which is a significant occupational health risk for plant operators and tunnel workers. The COSHH Regulations 2002 require that exposures to all these agents are assessed and controlled, that health surveillance is provided where WELs are likely to be exceeded, and that the results of health surveillance are used to monitor the effectiveness of control measures.}

\paragraph{Dermatitis and skin disease represent a further occupational health burden in construction that is frequently underestimated. Wet cement, epoxy resins, solvents, and cutting oils are among the most common causes of occupational contact dermatitis, which accounts for a significant number of construction workers each year being unable to work or requiring long-term treatment. Chromate sensitisation from wet cement causes allergic contact dermatitis that, once established, is permanent and can end a construction career. The Chromate in Cement Regulations now limit chromate content in cement, but the hazard has not been eliminated and wet cement skin contact remains a significant risk requiring gloves, barrier cream, and a change-out procedure.}

\barquote{``Silica dust does not announce itself. There is no pain, no immediate warning, and no recovery once the damage is done. The roofer cutting 50 blocks a day without a water-suppression saw is not aware that they are shortening their life. The employer who provides the dry angle grinder, and the site manager who watches the dust cloud without comment, are making a choice that will kill someone on a timescale that keeps them comfortable from accountability.''}

\example{Silicosis and Construction Site Enforcement}{An HSE inspector visits a commercial construction site during a phase of intensive brickwork and concrete blockwork. During the inspection, the inspector observes: multiple operatives using 4-inch angle grinders with cutting discs to cut bricks and blocks dry; visible dust clouds extending beyond the immediate work area; no water suppression in use; and no respiratory protective equipment (RPE) worn by any operative. Air monitoring conducted during the inspection produces RCS concentrations of 0.48 mg/m\textsuperscript{3} in the breathing zone of the most exposed operative --- almost five times the WEL. Improvement Notices are served requiring the immediate provision of on-tool water suppression for all block and brick cutting, the provision of appropriate RPE (minimum FFP3), and a COSHH assessment for the activities. A follow-up visit three months later finds partial compliance: some operatives are now using water-suppression attachments, but a number of operatives in a different work area are continuing to dry-cut without RPE. A Prohibition Notice is issued. Investigation reveals that the COSHH assessment, while recently prepared, had not been communicated to the bricklaying subcontractor's workforce and their foreman had not been briefed on the controls required. The systemic failure is the gap between the written assessment and the site-level implementation.}

%---------------------------- SECTION ----------------------------
\section{Legal and professional framework}

\paragraph{The primary regulatory framework for occupational health and dust in construction is the Control of Substances Hazardous to Health Regulations 2002 (COSHH), which require employers to assess the health risks arising from work with hazardous substances (including dusts, fumes, and biological agents), to prevent or adequately control exposure, and to provide health surveillance where appropriate. The COSHH framework is supplemented by the Workplace Exposure Limits published in EH40 (the HSE's document of occupational exposure limits), by specific Approved Codes of Practice for particular substances (including L5 for COSHH generally and specific ACOPs for woodworking and occupational asthma), and by a substantial body of HSE guidance documents and research papers. Key frameworks relevant to this chapter include: Health and Safety at Work etc.\ Act 1974; Management of Health and Safety at Work Regulations 1999; Control of Substances Hazardous to Health Regulations 2002; EH40 Workplace Exposure Limits (fourth edition 2020); Control of Asbestos Regulations 2012; Manual Handling Operations Regulations 1992; RIDDOR 2013; HSE guidance documents including EH44 (Dusty work --- controlling dust in construction); RR689 (Silica and construction); and the HSE COSHH essentials construction tool.}

\begin{itemize}
  \bitem{Control of Substances Hazardous to Health Regulations 2002 (COSHH)}{Require employers to: make a suitable and sufficient assessment of the risk to health from hazardous substances to which employees may be exposed; prevent or, where prevention is not reasonably practicable, adequately control that exposure; ensure that control measures are properly used and maintained; monitor exposure where appropriate; carry out health surveillance where appropriate; provide employees with information, instruction, and training; and ensure that emergency arrangements are in place. COSHH Regulation~7 requires that the WELs published in EH40 are not exceeded, and that exposure is reduced to as low as is reasonably practicable below the WEL.}
  \bitem{EH40 --- Workplace Exposure Limits (HSE, fourth edition 2020)}{Lists the WELs for all substances for which limits have been set in Great Britain, expressed as 8-hour time-weighted averages and, where applicable, short-term exposure limits (15-minute TWAs). The WELs for the most significant construction dust hazards are: respirable crystalline silica (quartz) 0.1 mg/m\textsuperscript{3} (8h TWA); inhalable dust 10 mg/m\textsuperscript{3} (8h TWA); respirable dust 4 mg/m\textsuperscript{3} (8h TWA); wood dust (hardwood) 3 mg/m\textsuperscript{3}; wood dust (softwood) 5 mg/m\textsuperscript{3}. The 2020 edition reduced a number of WELs including for respirable silica, reflecting updated scientific evidence.}
  \bitem{Control of Asbestos Regulations 2012}{Impose specific requirements for the assessment, notification, and control of asbestos exposure in construction, including the requirement for an R\&D survey before demolition or refurbishment, and the licensing requirements for work with ACMs in circumstances meeting the licensable work criteria. Asbestos is addressed in detail in Chapter~\ref{chap:Asbestos Risk Assessment}.}
  \bitem{RIDDOR 2013}{Requires the reporting of occupational diseases including: occupational dermatitis where it is caused by work; occupational asthma; and cases of Hand-Arm Vibration Syndrome. Silicosis and other dust-related diseases are not reportable when diagnosed (because they typically arise from historic exposures, often with multiple previous employers), but the duty to report occupational asthma cases under RIDDOR creates a management feedback loop that should inform the COSHH assessment review.}
  \bitem{HSE COSHH Essentials Construction tool}{The HSE's web-based tool for generating COSHH control guidance sheets for common construction activities, based on the substance being used, the activity being carried out, and the control strategy required. The tool provides practical guidance on the selection of appropriate controls and RPE for specific dust-generating operations, and is a useful reference for completing the COSHH assessment.}
\end{itemize}

%---------------------------- SECTION ----------------------------
\section{Conducting the assessment using the five-step method}

\subsection{Step 1 --- Identify Hazards}
\paragraph{Occupational health and dust hazards in construction fall into five primary categories: respirable mineral dusts (silica, asbestos fibres, mineral wool); organic dusts (wood, grain, flour, biological matter); chemical vapours and fumes (welding fume, isocyanates, epoxy resins, solvents, diesel engine exhaust); physical agents (noise, vibration --- addressed in Chapter~\ref{chap:Asbestos Risk Assessment}); and dermal hazards (wet cement, epoxy resins, solvents, nickel, chromates). The COSHH assessment must identify each specific substance or agent to which workers are exposed, the activity that generates the exposure, the route of exposure (inhalation, skin contact, ingestion), the likely concentration or level of exposure, and the WEL or BMGV (biological monitoring guidance value) applicable to the substance. Construction activities with the highest dust generation potential include: abrasive wheel cutting of concrete, brick, block, and stone; disc grinding on concrete and stone; pneumatic chipping and breaking of concrete and masonry; sand blasting for surface preparation; sanding and grinding of painted surfaces; demolition activities generating composite dust; and site sweeping with a brush rather than a wet or vacuum method.}

\subsection{Step 2 --- Decide Who or What Is Affected}
\paragraph{Workers directly conducting dust-generating operations are the primary risk group. However, the mobile nature of construction sites means that workers in adjacent trades, banksmen, visitors, and supervisors standing nearby are often exposed to dust clouds generated by others. The COSHH assessment must identify all persons in the exposure zone, not just the operative directly generating the dust. Workers in enclosed or poorly ventilated spaces --- tunnels, basement car parks, enclosed plant rooms, confined access areas --- are at higher risk due to reduced dilution ventilation. Workers who are young, who are relatively new to the industry, and who will have a longer remaining working life are at higher cumulative risk from chronic occupational disease exposures, and this must be taken into account when prioritising the provision of health surveillance.}

\subsection{Step 3 --- Evaluate Likelihood and Consequence}
\paragraph{The consequence of chronic silicosis is incurable, progressive, and frequently fatal respiratory impairment and premature death. The consequence of occupational lung cancer from silica or diesel exposure is fatal. The consequence of occupational asthma is a permanent, career-ending condition. These outcomes are Catastrophic on any risk matrix, even though the probability of the consequence materialising in any individual exposure event is low. The cumulative nature of the hazard means that the risk assessment cannot treat each day's exposure independently: the assessment must evaluate the cumulative lifetime exposure risk, which requires knowledge of likely exposure concentrations and the duration of the role or career over which the worker will be exposed. A young worker entering the construction industry today and exposed to silica at twice the WEL for 30 years is virtually certain to develop silicosis.}

\subsection{Step 4 --- Select and Implement Controls}
\paragraph{The COSHH hierarchy of control requires prevention of exposure as the first priority, substitution of less hazardous substances or methods, engineering controls, administrative controls, and PPE in that order. For silica dust: the primary engineering control is wet cutting (water supplied at the blade to suppress dust at source) or on-tool extraction (local exhaust ventilation coupled to the cutting tool, capturing dust before it is released to the atmosphere). On-tool extraction requires a dust extraction unit with the appropriate filter efficiency (typically H-class) and a power tool fitted with a compatible extraction shroud. Wet cutting at the rate of water supply and pressure required to provide effective suppression must be used for angle grinder cutting where on-tool extraction is not available. Where neither wet cutting nor on-tool extraction fully suppresses exposure below the WEL, RPE (minimum FFP3 disposable filtering face piece) must be worn as an additional control, not as a substitute for engineering controls.}

\subsection{Step 5 --- Record and Review}
\paragraph{The COSHH assessment for dust-generating activities must be documented and communicated to all operatives and sub-contractors who carry out the activities on site. Health surveillance records, including spirometry results and skin surveillance records, must be maintained in accordance with the COSHH Regulations for the period prescribed (40 years for asbestos health records; 5 years for other COSHH health records). Air monitoring results must be retained and must be used to verify that engineering controls are achieving the required reduction in exposure. The COSHH assessment must be reviewed whenever work methods change, new substances or activities are introduced, health surveillance data indicates that exposures may not be as controlled as assumed, or a new WEL comes into force.}

\example{Five-Step Application}{A groundworks and drainage contractor is constructing manholes in reinforced concrete using a vibrating poker, disc saw for slab trimming, and angle grinders for rebar cutting. \textbf{Step~1 (Hazards):} Respirable crystalline silica from concrete disc saw cutting; inhalable cement dust from mixing and placing; concrete dust from angle grinder operations on cured concrete. \textbf{Step~2 (Who Affected):} Concretor, banksman in vicinity, drainage operatives working adjacent. \textbf{Step~3 (Evaluation):} Silica rated Catastrophic (long-term)/Medium probability of significant exposure without controls. \textbf{Step~4 (Controls):} Wet attachment on disc saw providing continuous water suppression at blade; on-tool extraction shroud on angle grinder connected to H-class extractor; FFP3 RPE worn by operator and any person within 2 m; dry sweeping of concrete dust prohibited --- vacuum or wet method only; daily inspection of water suppression supply before cutting commences. \textbf{Step~5 (Record):} COSHH assessment communicated to sub-contractor supervisor at TBT; quarterly air monitoring for RCS to verify controls; health surveillance (spirometry) offered to concretors at annual intervals.}

%---------------------------- SECTION ----------------------------
\section{Applying the hierarchy of controls}

\paragraph{The COSHH hierarchy, as required by COSHH Regulation~7, places prevention of exposure at the apex and RPE at the base. On construction sites, there is a persistent tendency to skip directly to RPE, which is the least reliable and most dependent on human behaviour of all the control measures.}

\begin{itemize}
  \bitem{Elimination}{Eliminate the dust-generating activity entirely where possible. Pre-cut concrete panels, pre-formed block sizes, and pre-cut timber eliminate the need for on-site cutting. Specify materials that do not require dust-generating machining at the point of use. Pre-mixed mortars delivered in ready-to-use sealed units eliminate the need for dry cement mixing.}
  \bitem{Substitution}{Substitute high-silica-content materials with lower-silica alternatives where technically acceptable. Substituting engineered stone or granite composite worktops with lower-silica content materials for kitchen and bathroom fitting has been the subject of specific HSE guidance following a spike in rapid-onset silicosis (silico-proteinosis) in stone fabricators. Substitute dry cutting with wet cutting. Substitute open-site concrete mixing with pre-mix or premix-in-bag systems.}
  \bitem{Engineering controls}{On-tool extraction (local exhaust ventilation coupled to the cutting tool, capturing dust at source) and wet cutting with adequate water supply are the primary engineering controls for silica dust. General ventilation by natural or mechanical means reduces background dust levels in enclosed work areas. Damping down of dusty surfaces and haul roads reduces resuspension of settled dust. On dusty demolition sites, water misting systems can reduce ambient dust concentrations across the working area.}
  \bitem{Administrative controls}{Permit-to-work systems for dust-generating activities in enclosed spaces. COSHH assessment communicated to all operatives and sub-contractors before work commences. Prohibition on dry sweeping of concrete or silica-containing dust. Rotation of workers to reduce individual cumulative exposure duration. Health surveillance programme for workers regularly exposed to silica, wood dust, or sensitising chemicals.}
  \bitem{PPE}{Where engineering controls do not reduce exposure below the WEL, RPE of the appropriate assigned protection factor (APF) must be provided. For RCS exposures up to ten times the WEL, an FFP3 disposable filtering face piece (APF 20) is the minimum. For higher exposures, a half-mask with P3 filters (APF 20), a powered air-purifying respirator (PAPR) with TH3 hood (APF 20), or a full-face mask with P3 filters (APF 40) may be required. RPE must be fit-tested (for tight-fitting face pieces), inspected before each use, and maintained in accordance with the manufacturer's instructions. Impermeable gloves and barrier cream for wet cement and epoxy resin contact.}
\end{itemize}

%---------------------------- SECTION ----------------------------
\section{Cadence of assessment and review triggers}

\paragraph{The COSHH assessment for dust-generating activities must be reviewed at defined intervals and following specific trigger events.}

\begin{itemize}
  \bitem{Before commencement of each new dust-generating activity}{Each distinct dust-generating activity --- a new cutting operation, a new phase of block-work, a new concrete pour involving cutting --- must be assessed before it commences. If the activity is covered by an existing COSHH assessment that has been reviewed and is current, the assessment need not be redone, but its application to the specific activity must be confirmed.}
  \bitem{When health surveillance data indicates poor control}{If spirometry results for workers regularly exposed to silica show declining lung function, or if occupational asthma is diagnosed or reported in a worker on the site, the COSHH assessment for the relevant activities must be reviewed immediately, air monitoring conducted, and controls improved. Health surveillance data is the feedback mechanism for COSHH compliance; it must not simply be filed without review.}
  \bitem{When air monitoring results exceed the WEL}{Air monitoring is not required for all COSHH activities but is best practice for silica-generating activities where exposure levels are uncertain. If monitoring results show exposure above the WEL, the assessment must be reviewed and controls enhanced until monitoring confirms compliance. Monitoring results below the WEL confirm that controls are effective and the assessment remains valid.}
  \bitem{When new materials or activities are introduced}{Any new material introduced to the site that contains a hazardous substance, or any new activity that generates dust, fume, or vapour, must be assessed before use. The safety data sheet (SDS) for the new material must be obtained and reviewed.}
  \bitem{At the frequency required by COSHH Regulation 11 for health surveillance}{Health surveillance for silica exposure must be provided to workers at risk at intervals not exceeding the COSHH Approved Code of Practice recommendations, typically annually for spirometry where regular silica exposure occurs.}
\end{itemize}

%---------------------------- SECTION ----------------------------
\section{Common failure modes}

\begin{itemize}
  \bitem{COSHH assessment prepared but controls not communicated to operatives}{The most common COSHH compliance failure on construction sites is the existence of a written assessment that correctly identifies the controls required but has never been communicated to the operatives or sub-contractors who carry out the work. The assessment specifies wet cutting; the operative uses a dry angle grinder because nobody told them otherwise. Effective COSHH compliance requires that the controls are briefed at toolbox talk, reinforced by supervision, and verified by periodic checking.}
  \bitem{RPE used without fit testing}{Tight-fitting disposable FFP3 face pieces are frequently worn on construction sites by operatives who have never been fit-tested, by operatives with beards that prevent an adequate seal, or over glasses frames that break the seal at the side of the face. An improperly fitting FFP3 provides substantially less than its rated protection. Fit testing of all tight-fitting RPE is required before first use, and must be repeated if the wearer's facial features change.}
  \bitem{On-tool extraction used but not maintained}{Local exhaust ventilation systems coupled to cutting tools lose their effectiveness rapidly if filters are not replaced, if filter bags become full, or if the extraction hose develops a leak. An H-class dust extractor with a blocked filter is effectively no control at all. The maintenance of on-tool extraction equipment must be documented, and daily pre-use checks must include filter condition and motor performance.}
  \bitem{Dry sweeping of concrete and silica-containing dust}{The prohibition on dry sweeping of silica-containing dust is specified in virtually every COSHH assessment for concrete work but continues to be observed on construction sites. Dry sweeping of hardened concrete dust generates RCS concentrations that can exceed the WEL by a factor of ten or more. Vacuum sweeping or wet methods must be specified and enforced.}
  \bitem{No health surveillance programme for at-risk workers}{Employers who regularly expose workers to silica, hardwood dust, isocyanates, or other COSHH-listed substances requiring health surveillance frequently have no health surveillance programme in place. The absence of health surveillance means that early-stage occupational disease --- the stage at which intervention could still protect the worker's health --- goes undetected until irreversible damage has occurred.}
\end{itemize}

\example{Failure Pattern}{A domestic extension contractor employs three site carpenters as full-time employees. The carpenters carry out all site woodworking including roof truss cutting, window and door frame preparation, and floor boarding, primarily using circular saws and routers without on-tool extraction. No COSHH assessment for wood dust has ever been prepared. Three years after the company has been trading, one carpenter, aged 44, presents to his GP with persistent nasal congestion that has been worsening for two years. An ENT investigation identifies a sinonasal malignancy. The worker is referred to an occupational physician who attributes the malignancy, on the balance of probabilities, to occupational hardwood dust exposure. A civil claim is brought against the employer. The employer's insurance company discovers that no COSHH assessment was ever prepared, no wood dust controls were in place beyond occasional dust masks of unknown specification, and no health surveillance had been offered. The claim is settled for a substantial sum.}

%---------------------------- CHAPTER SUMMARY ----------------------------
\newpage
\section{Chapter Summary}

\begin{table}[H]
\centering
\begin{tabularx}{\textwidth}{>{\raggedright\arraybackslash}p{3.2cm}X}
\toprule
\textbf{Assessment Element} & \textbf{Operational Requirement} \\
\midrule
Legal/Professional Basis & COSHH Regulations 2002; EH40 WELs; CAR 2012; RIDDOR 2013; Health and Safety at Work etc.\ Act 1974; Management Regulations 1999. \\
Method & Apply the full COSHH five-step assessment for each dust and chemical hazard. Identify the substance, the activity, the exposure route, the WEL, and the controls required to reduce exposure to as low as reasonably practicable below the WEL. \\
Controls & On-tool extraction or wet cutting as primary engineering controls for silica dust. RPE (minimum FFP3 fit-tested) as supplementary control only. Prohibition on dry sweeping. Health surveillance for at-risk workers. Communication of COSHH assessment to all operatives and sub-contractors. \\
Cadence & Before each new dust-generating activity; when health surveillance data indicates poor control; when air monitoring exceeds WEL; when new materials introduced; at health surveillance intervals. \\
Assurance & Air monitoring records; health surveillance records (retained 40 years for asbestos, 5 years for other COSHH); RPE fit-test records; on-tool extraction maintenance log; COSHH assessment review records. \\
\bottomrule
\end{tabularx}
\end{table}

\begin{itemize}
  \bitem{Key takeaway 1}{Occupational disease kills far more construction workers each year than acute traumatic injury. Silicosis, occupational lung cancer, mesothelioma, and occupational asthma are preventable through COSHH assessment and control. The invisibility of the mechanism does not reduce the employer's legal duty or moral obligation.}
  \bitem{Key takeaway 2}{On-tool extraction and wet cutting are the primary engineering controls for silica dust generated by cutting operations. RPE is a supplementary control, not a substitute for engineering controls. An operative cutting blocks dry while wearing an FFP3 face piece is partially protected; the same operative cutting wet or with on-tool extraction is properly protected.}
  \bitem{Key takeaway 3}{RPE must be fit-tested before first use, inspected before each use, and maintained in accordance with the manufacturer's instructions. Beards, spectacles, and facial features that prevent an adequate seal make tight-fitting respirators ineffective regardless of their filter efficiency rating.}
  \bitem{Key takeaway 4}{Health surveillance for workers regularly exposed to silica, hardwood dust, and sensitising chemical agents must be provided at the intervals required by the COSHH Regulations. Spirometry results showing declining lung function must trigger an immediate review of exposure controls.}
  \bitem{Key takeaway 5}{The COSHH assessment must be communicated to all operatives and sub-contractors before dust-generating activities commence. A written assessment that has never been briefed to the people it is designed to protect is not an effective control.}
\end{itemize}

\barquote{In Chapter~\ref{chap:Welfare, Site Facilities and Vulnerable Workers Risk Assessment}, we move into Part~6, examining the legal requirements and practical arrangements for welfare facilities on construction sites --- the foundation of a site safety culture and the indicator, above all others, of the value that an organisation places on the people who do the work.}
